\section{经典数据集：BERT、WebText、CCNet、C4}

BERT 使用 Wikipedia 和 BooksCorpus。BooksCorpus 来自 Smashwords 上免费的自出版书籍，后来因违反 ToS 等问题下线。GPT-2 的 WebText 则用 Reddit 高 karma 外链作为 quality surrogate，并产生 OpenWebText 的开源复刻。

\begin{knowledgebox}{课堂提示：Reddit karma 只是 quality surrogate}
WebText 没有直接判断网页是否“高质量”，而是收集 Reddit 帖子中 karma 不低于 3 的外链页面，把社区投票当作代理信号。老师把它作为一个典型启发式：代理信号便宜、可扩展，却同时继承社区人口、兴趣和平台机制的偏差，不能被误读成客观质量标签。
\end{knowledgebox}

CCNet 试图自动构造大规模高质量预训练数据：去重、语言识别、用 KenLM 5-gram 模型筛出像 Wikipedia 的文本。C4 从 Common Crawl 2019 snapshot 出发，用手工规则过滤自然语言文本，成为 T5 的重要贡献之一。

\begin{figure}[H]
\centering
\includegraphics[width=0.88\textwidth]{c4-domains.png}
\caption{C4 domain analysis：规则过滤后的 Common Crawl 仍有明显 domain 组成和偏差。}
\figsource{Stanford CS324 lecture image，本地化后供 CS336 Lecture 13 使用。}
\end{figure}

\begin{knowledgebox}{读图：C4 domains 告诉我们什么}
C4 不只是“网页文本”。它的 domain distribution 反映了 crawl、HTML extraction 和过滤规则共同造成的选择偏差。模型学到的风格、知识和价值取向会受这些 domain 比例影响。
\end{knowledgebox}

\begin{importantbox}{经典数据集的谱系}
\begin{longtable}{L{0.13\textwidth}L{0.23\textwidth}L{0.23\textwidth}L{0.17\textwidth}}
\toprule
数据集 & 主要来源 & 关键处理 & 教学结论 \\
\midrule
\endhead
BERT data & Wikipedia + BooksCorpus & 文档级序列，BooksCorpus 来自免费自出版书。 & 早期 LM 数据更窄、更人工。 \\
WebText & Reddit 高 karma 外链网页 & 英文过滤、近重复移除。 & 用社交信号做代理质量筛选。 \\
CCNet & Common Crawl & 去重、fastText 语言识别、KenLM 质量过滤。 & 自动化 web filtering 的早期范式。 \\
C4 & Common Crawl snapshot & 手工规则、坏词/非英文/短页过滤。 & 简单规则可规模化，但偏差明显。 \\
GPT-3 data & CC、WebText2、Books、Wikipedia & 质量分类器、fuzzy dedup。 & 私有数据细节开始变模糊。 \\
\bottomrule
\end{longtable}
\end{importantbox}

